% !TeX TS-program = xelatex

\documentclass{resume}
\ResumeName{李树臻}

% 如果想插入照片，请使用以下两个库。
% \usepackage{graphicx}
% \usepackage{tikz}

\begin{document}
\linespread{1.25}

\ResumeContacts{
  (+86)132-8778-1587,%
  \ResumeUrl{mailto:}{2538152560@qq.com},%
}

% 如果想插入照片，请取消此代码的注释。
% 但是默认不推荐插入照片，因为这不是简历的重点。
% 如果默认的照片插入格式不能满足你的需求，你可以尝试调整照片的大小，或者使用其他的插入照片的方法。
% 不然，也可以先渲染 PDF 简历，然后用其他工具在 PDF 上叠加照片。
% \begin{tikzpicture}[remember picture, overlay]
%   \node [anchor=north east, inner sep=1cm]  at (current page.north east) 
%      {\includegraphics[width=2cm]{image.png}};
% \end{tikzpicture}

\ResumeTitle


\section{教育经历}
\ResumeItem
[深圳大学|硕士研究生]
{深圳大学}
[\textnormal{集成电路工程，电子与信息工程学院|}  专业型硕士研究生]
[2024.09—2027.06（预计）]

\textbf
主要研究方向为深度学习算法，在计算机视觉领域方面有一定的研究和工程经验。\textbf{2027年应届生}。

主要研究成果为： Signal Processing: Image Communication 期刊的论文《SWALIC: Asymmetric Wavelet Image Compression
with Parallel Shift-Attention Architecture and
Cross-Band Collaborative Refinement》。 (在投)

\ResumeItem
[曲阜师范大学|本科生]
{曲阜师范大学}
[\textnormal{计算机科学与技术，计算机学院|} 工学学士]
[2020.09—2024.06]

\textbf{GPA: 4.25/5.0(专业前 1\%)}，获学业奖学金多次，中国高校计算机大赛-团体程序设计天梯赛 全国团队二等奖，第十二届山东省ICPC大学生程序设计竞赛 铜奖。

\section[技术能力]{技术能力}
\begin{itemize}
  \item \textbf{语言}: Python, C, C++, Verilog等
  \item \textbf{算法}：掌握程序设计竞赛的相关算法，对贪心，搜索和图论等较为了解
  \item \textbf{英语}：CET-4, CET-6
\end{itemize}

\section{项目经历}

\ResumeItem{\textbf{AELIC-Wave} 非对称轻量学习图像压缩}
[深度学习 / 图像压缩]
[2025.03—2026.12]
\begin{itemize}
  \item 提出面向\textbf{资源受限边缘设备（如无人机、遥感卫星等）}的基于两层 Bior4.4 小波分解的子带感知非对称 VAE 框架，实现编码端极低复杂度。
  \item 设计 Cross-Band Modulation Module (CBMM) 利用 LL 子带先验预增强 HF latent；设计 Contextual Refinement Module (CRM) 和 PostProcess Module 实现解码端渐进精炼。
  \item 通过完备消融实验证明 CBMM 需全套解码器协同才能发挥全部增益。
  \item 在 AID/UCM 遥感数据集上完成训练与评估，R-D 性能与先进轻量方法可比，计算开销降低超 90\%，编码器仅 0.166M 参数和 0.063GMACs。
\end{itemize}

\ResumeItem{\textbf{SWALIC} 非对称轻量学习图像压缩}
[深度学习 / 图像压缩]
[2026.01—2026.06]
\begin{itemize}
  \item 面向无人机及遥感卫星等嵌入式边缘平台，在 AELIC-Wave 框架基础上，将卷积骨干全面替换为并行注意力-移位骨干（LSS/SBB/PASB）。
  \item 设计可学习空间移位（LSS），实现子带特异的高效特征混合；设计 SBB（带扩展比的 shift bottleneck）和 PASB（并行门控注意-移位融合块）。
  \item 在 AID/UCM 遥感数据集上完成训练与评估，R-D 性能在AELIC-Wave 框架基础上进一步提升，并评估了 MACs、CPU 延迟、峰值内存等多个维度，计算开销在AELIC-Wave 框架基础上进一步降低，编码器仅 0.119M 参数和 0.044GMACs。
\end{itemize}

\ResumeItem{\textbf{DVS} 手势识别 — 轻量级点云时序网络}
[深度学习 / 事件相机]
[2025.04—2026.03]
\begin{itemize}
  \item 针对 DVS 事件流稀疏含噪的问题，设计并实现了一套轻量级手势识别系统，包含时空网格去噪、点云空间特征提取与 TCN时序特征提取三个模块。
  \item 在 DVS128 Gesture 数据集（11 类手势）上以 \textbf{68K 参数}达到 \textbf{95.12\%} 准确率，单次推理仅需 19 MFLOPs、768KB 内存，适用于资源受限的边缘设备。
  \item 搭建实时流式识别系统，结合滑动窗口与投票滤波，在 50ms 延迟下将在线准确率提升至 \textbf{97.1\%}，适用于机器人交互、智能穿戴等资源受限场景。
\end{itemize}




\end{document}
